Nós dizemos que um subconjunto de \(\mathbb{R}^n\) é \(k\)-\textit{quase contido} por um hiperplano se há menos de \(k\) pontos nesse conjunto que não pertencem ao hiperplano. Chamamos um conjunto finito de pontos de \(k\)-\textit{genérico} se não existe hiperplano que \(k\)-quase contém o conjunto. Para cada par de inteiros positivos \((k, n)\), encontre o número mínimo de \(d(k, n)\) tal que todo conjunto finito \(k\)-genérico em \(\mathbb{R}^n\) contém um subconjunto \(k\)-genérico com no máximo \(d(k, n)\) elementos.
